\section{Products and the Evaluation Crisis}
\label{sec:products}
\label{sec:eval}

A small industry now sells the folkloric dream as a feature. The archetype
is Takara's \sys{Bow-Lingual} (2002): a collar microphone and handheld
display that matched bark acoustics to six emotions and then to one of
roughly 200 canned phrases~\cite{bowlingual2002}. It sold well, won the 2002
Ig Nobel Peace Prize, and established the template: classify affect, then
decorate it with a sentence. Two decades later the template is unchanged;
only the classifier has been swapped for a deep model and the phrase bank
for a language model. The current flagship, \sys{Traini}, analyzes sound,
images, and video to infer emotion and need, reporting about 12 emotions
across roughly 120 breeds with a headline accuracy of
81.5\%~\cite{traini2024}. Its founder has publicly conceded that audio alone
is insufficient and that body language and face must be fused to reach the
claimed accuracy, an implicit admission that the product is multimodal
affect inference; the reverse direction (human speech rendered as ``dog'')
is far weaker, on the order of 40\% in independent coverage.
\sys{MeowTalk} applies the idea to cats~\cite{meowtalk2020}, the
\sys{Petpuls} collar classifies emotional states from
vocalizations~\cite{petpuls2021}, and in China ``Meng Xiao Yi'' reports a
Qwen-based model trained on millions of pet voiceprints with a claimed
94.6\% emotion accuracy~\cite{mengxiaoyi2026}. Table~\ref{tab:products}
audits the landscape against the ladder.

Read carefully, every headline number is an emotion-classification metric,
not a translation metric. ``81.5\%'' means that on some held-out set the
model's emotion label agreed with a human-assigned label 81.5\% of the
time. That is a real, falsifiable Rung-2 result, if the test set, label
protocol, and class balance are disclosed, which consumer vendors rarely
do. It says nothing about whether the rendered sentence corresponds to
anything the dog intended, because no such ground truth exists. The
sentence is generated downstream of the classifier by a language model
optimizing for human emotional resonance, not fidelity to the animal. The
product demonstrates Rung 2 and sells Rung 5.

Three incentives keep the gap open. The consumer value proposition is
emotional, so a plausible affectionate sentence beats a calibrated
probability vector even though the latter is the honest output. The data
flywheel favors engagement over validation: every uploaded clip enriches
the classifier, but no mechanism tests whether the ``translation'' is
correct, because the dog cannot complain. And the scientific validation
loop is slow and unglamorous and yields no shareable screenshot. The result
is a market optimized for the appearance of translation.

The deeper problem is falsifiability. Translation in the human sense is
scored against a reference; canine communication has none. There is no
canonical intended message behind a bark and no bilingual informant who can
confirm that ``I want to play'' is what the dog meant, and per
Section~\ref{sec:biology} there may be no sentence-shaped intent to report.
What can be scored is the upstream classifier, and reporting that as a
``translation accuracy'' is a category mistake. The ELTE tradition is
explicit that its labels are human-perceived emotion or eliciting
context~\cite{pongracz2006acoustic,icassp2018dogbark}. ``This bark means `I
love you'\,'' admits no observation that could show it false: the dog will
not contradict it, and any subsequent behavior can be retrofitted to the
narrative. Compare a weather forecast, scored against tomorrow's weather,
or speech-to-text, scored against a transcript. The pet translator has no
tomorrow and no transcript.

Two classic threats compound this. The Clever Hans effect, the horse that
appeared to do arithmetic but was reading its trainer's unconscious cues,
is the permanent hazard of animal-communication experiments; soundboard
studies guard against it with blind protocols and unfamiliar
button-pressers~\cite{bastos2024plos}, while consumer apps used by the
devoted owner in an uncontrolled home do the opposite. Anthropomorphism is
the second: human observers project their categories onto the
dog~\cite{sci2022explainable}, and a language model trained on human text
amplifies the projection, risking teaching owners to ignore the dog's real
body signals in favor of a flattering sentence.

A defensible evaluation program therefore has four ingredients: defined
labels and disclosed protocols on a public held-out corpus, the BEANS/BirdSet
discipline applied to dogs~\cite{hagiwara2023beans}; calibration, so
probabilities match observed frequencies and a welfare tool can threshold
responsibly; closed-loop validation of inferred intent against the dog's
subsequent behavior under blind conditions; and a strict separation of the
rendered sentence from the score. None of today's consumer products
satisfies these. Table~\ref{tab:evidence} maps the principal canine results
to their actual label source: everything at Rungs~1--2 has an objective or
consensus label and is falsifiable; Rung~3 is falsifiable only because a
human-designed symbol set makes the intent checkable; the bark-to-sentence
claim has no label at all.

This crisis is the animal-communication instance of a broader reckoning in
AI about what benchmarks measure. The ``stochastic parrots'' critique warns
that scaling text models amplifies form without
grounding~\cite{bender2021parrots}, the ``climbing toward NLU'' position
argues that the field mistakes benchmark gains for
understanding~\cite{bender2020climbing}, and broad suites (HELM, BIG-bench,
MMLU, TruthfulQA) expose how much headline scores depend on task choice and
contamination~\cite{liang2023helm,srivastava2023bigbench,hendrycks2021mmlu,lin2022truthfulqa}.
Documentation practices such as datasheets and model cards exist to make a
system's limits legible~\cite{gebru2021datasheets,mitchell2019modelcards},
and the hallucination literature shows that fluent generation routinely
outruns its evidence~\cite{ji2023hallucination,maynez2020faithfulness}.
Table~\ref{tab:limits} consolidates the limitations this survey has
identified. We propose a simple demarcation: a canine communication system
should be called a translator only if its outputs are scored against an
independent, animal-grounded error signal. By this criterion no current
product qualifies; the soundboard paradigm qualifies at the word level; the
rest is affect classification with a narrative interface. This is the
difference between a measurement and a story.

\begin{table}[htbp]
\centering
\caption{Commercial ``pet translators'' audited against the ladder.
Accuracies are vendor-reported claims, not independently verified; all
measure emotion classification (Rung 2).}
\label{tab:products}
\footnotesize
\setlength{\tabcolsep}{4pt}
\begin{tabular}{@{}p{2.5cm}p{3.5cm}p{2.6cm}p{2.0cm}p{2.2cm}@{}}
\toprule
\textbf{Product} & \textbf{Modality / method} & \textbf{Claim} & \textbf{Actual rung} & \textbf{Reverse} \\
\midrule
Bow-Lingual (2002)~\cite{bowlingual2002} & bark acoustics $\to$ 6 emotions & $\sim$200 phrases & 2 & no \\
Traini / PetGPT~\cite{traini2024} & audio + image + video & 81.5\%, 12 emotions & 2 & weak ($\sim$40\%) \\
MeowTalk~\cite{meowtalk2020} & cat vocalization classifier & emotion labels & 2 & no \\
Petpuls~\cite{petpuls2021} & collar vocalization analysis & emotion states & 2 & no \\
Meng Xiao Yi~\cite{mengxiaoyi2026} & Qwen + voiceprints & 94.6\% (claimed) & 2 & no \\
DogTok & human command $\to$ bark & training aid & n/a & yes \\
\bottomrule
\end{tabular}
\end{table}

\begin{table}[htbp]
\centering
\caption{Evidence-type map: what ground truth each canine study actually
has. ``Falsifiable'' means an independent error signal exists.}
\label{tab:evidence}
\footnotesize
\setlength{\tabcolsep}{4pt}
\begin{tabular}{@{}p{3.0cm}p{3.4cm}p{3.3cm}p{1.7cm}p{2.2cm}@{}}
\toprule
\textbf{Study} & \textbf{Task} & \textbf{Label / ground truth} & \textbf{Rung} & \textbf{Falsifiable?} \\
\midrule
Pongr\'acz et al.~\cite{pongracz2005barks,pongracz2006acoustic} & bark context / emotion & human consensus on context & 2 & yes \\
Moln\'ar et al.~\cite{molnar2006discriminate} & individual ID & dog identity (objective) & 1 & yes \\
Farag\'o et al.~\cite{farago2017growls} & growl size / intent & human judgment + known context & 2--3 & yes \\
ICASSP 2018~\cite{icassp2018dogbark} & context / emotion & human-perceived label & 2 & yes, if disclosed \\
DogFACS CNNs~\cite{sci2022explainable,sci2023cnnewoa,sci2025automated} & facial affect / AUs & human-coded FACS & 2 & yes \\
Bastos et al.\ (PLoS)~\cite{bastos2024plos} & word comprehension & blind behavioral response & 3 & yes (scaffolded) \\
Bastos et al.\ (Sci.\ Rep.)~\cite{bastos2024scirep} & two-button combination & combination statistics & 3 & yes (scaffolded) \\
bark-to-sentence apps~\cite{traini2024} & rendered sentence & none (no intended message) & 5 (claimed) & \textbf{no} \\
\bottomrule
\end{tabular}
\end{table}

Figure~\ref{fig:evidence} reads the two audits together. Counting the entries
of Tables~\ref{tab:evidence} and \ref{tab:products} by the rung each one
actually reaches gives a shape no single table shows: seven falsifiable
canine studies cluster at rungs 1--3, five consumer products cluster at
rung 2, and right of the gap the validated column is empty while the claim
column holds one entry --- the bark-to-sentence category, scored at the rung
it advertises. Evidence and marketing here differ less in quality than in
location.

\begin{figure}[htbp]
\centering
\resizebox{0.94\textwidth}{!}{% Figure - Evidence vs. claims along the ladder.
% Left series: rows of Table tab:evidence, i.e. canine studies with an
% independent error signal, assigned to the rung each study actually reaches
% (the growl study spans rungs 2--3 and is counted at rung 2).
% Right series: rows of Table tab:products, i.e. audited consumer products,
% assigned to their actual rung; the single Rung-5 entry is the
% bark-to-sentence product category, whose claim is unfalsifiable.
\begin{tikzpicture}[font=\sffamily]
  \begin{axis}[
      name=ev,
      width=13.6cm, height=5.9cm,
      ymin=0, ymax=7.4,
      symbolic x coords={1,2,3,4,5},
      xtick=data,
      xticklabels={%
        {\color{skydeep}\bfseries 1 \textperiodcentered\ Signal},
        {\color{sagedeep}\bfseries 2 \textperiodcentered\ Affect},
        {\color{gold!60!black}\bfseries 3 \textperiodcentered\ Intent},
        {\color{pawdeep}\bfseries 4 \textperiodcentered\ Reference},
        {\color{rosedeep}\bfseries 5 \textperiodcentered\ Dialogue}},
      ytick={0,1,2,3,4,5,6},
      tick label style={font=\scriptsize, text=ink!70},
      axis lines*=left, line width=0.6pt, draw=ink!45,
      ylabel near ticks, ylabel style={font=\sffamily\scriptsize, text=ink!70},
      ylabel={catalogued entries},
      grid=major, grid style={ink!8},
      ybar=7pt, bar width=19pt,
      every node near coord/.style={font=\sffamily\bfseries\scriptsize,
        text=ink, anchor=south, inner sep=1.5pt},
      enlarge x limits=0.09,
    ]
    % shaded region: above the gap, nothing has been validated
    \path (axis cs:3,0) -- (axis cs:4,0) coordinate[midway] (gapx);
    \fill[gap, opacity=0.45] (gapx |- current axis.south)
      rectangle (current axis.north east);
    \addplot+[ybar, mark=none, fill=sagedeep!85, draw=sagedeep, line width=0.3pt,
              nodes near coords]
      coordinates {(1,1) (2,4) (3,2) (4,0) (5,0)};
    \addplot+[ybar, mark=none, fill=rosedeep!70, draw=rosedeep, line width=0.3pt,
              nodes near coords]
      coordinates {(1,0) (2,5) (3,0) (4,0) (5,1)};
    \draw[dashed, rosedeep, line width=1.1pt] (gapx |- current axis.south)
      -- (gapx |- current axis.north);
  \end{axis}

  % hand-drawn legend: solid swatches in one row, clear of the bars
  \node[fill=sagedeep!85, draw=sagedeep, line width=0.3pt, rounded corners=1.5pt,
        minimum width=0.5cm, minimum height=0.26cm, inner sep=0, anchor=west] (sw1)
    at ($(ev.north west)+(0.30cm,-0.42cm)$) {};
  \node[font=\sffamily\scriptsize, text=ink, anchor=west] (l1)
    at (sw1.east) {\;falsifiable canine studies};
  \node[fill=rosedeep!70, draw=rosedeep, line width=0.3pt, rounded corners=1.5pt,
        minimum width=0.5cm, minimum height=0.26cm, inner sep=0, anchor=west] (sw2)
    at ($(l1.east)+(0.55cm,0)$) {};
  \node[font=\sffamily\scriptsize, text=ink, anchor=west] (l2)
    at (sw2.east) {\;audited consumer products};

  \node[font=\sffamily\bfseries\scriptsize, text=rosedeep, align=left,
        anchor=north west]
    at ($(gapx |- ev.north west)+(1.9cm,-2pt)$)
    {THE SEMANTIC GAP\\[1pt]{\sffamily\tiny\normalfont zero validated entries to the right}};
  \node[font=\sffamily\bfseries\scriptsize, text=ink!75, anchor=south west]
    at ($(ev.north west)+(0pt,4pt)$) {evidence thickens below the gap, claims pile up above it};
\end{tikzpicture}
}
\caption{Catalogued entries per rung of the ladder. Green: the canine studies
of Table~\ref{tab:evidence} that carry an independent error signal, assigned
to the rung each study reaches (the growl study spans rungs 2--3 and is
counted at rung 2; the bark-to-sentence row is excluded from this series).
Rose: the audited products of Table~\ref{tab:products} at their actual rung
(DogTok, a command-to-bark aid, has no receptive rung and is left out), plus
the bark-to-sentence category placed at the rung it claims. The shaded
region is the semantic gap of Figure~\ref{fig:ladder}: nothing to its right
has been validated.}
\label{fig:evidence}
\end{figure}

\begin{table}[htbp]
\centering
\caption{Consolidated limitations of canine ``translation,'' with locus and
consequence.}
\label{tab:limits}
\footnotesize
\setlength{\tabcolsep}{4pt}
\begin{tabular}{@{}p{4.0cm}p{4.6cm}p{5.0cm}@{}}
\toprule
\textbf{Limitation} & \textbf{Where it bites} & \textbf{Consequence} \\
\midrule
No ground-truth message & bark-to-sentence scoring & fidelity is unfalsifiable \\
Human-perceived labels & Rung-2 metrics & a proxy for state, not a readout \\
Undisclosed protocols & vendor accuracies & claims cannot be checked \\
Clever Hans / anthropomorphism & uncontrolled home use & the observer leaks the answer \\
No parallel corpus & supervised training & no translation signal exists \\
Semantic gap / no recursion & Rung 5 & the target may not exist \\
Fluency $\neq$ fidelity & LLM sentence rendering & hallucinated ``translations'' \\
No closed-loop error signal & consumer market & the field cannot self-correct \\
\bottomrule
\end{tabular}
\end{table}
